\chapter{Grundlagen der Sensorik und Hardware}

Dieses Kapitel behandelt die grundlegenden Konzepte der Sensorik und die für dieses Projekt verwendete Hardware. Zunächst werden die allgemeinen Prinzipien der Sensorik erläutert sowie deren Relevanz für die Bewegungserkennung. Anschließend folgt eine detaillierte Vorstellung der eingesetzten Hardware-Komponenten. Abschließend werden die verwendeten Sensoren beschrieben und ihre Bedeutung für die Erfassung und Analyse von Bewegungen dargelegt.

\section{Einführung}
Zu Start, gibt es eine kurze Einführung, Was Sensorik und Hardware sind, sowie inwiefern diese relevant ist für die Bewegungserkennung
\subsection{Was ist Sensorik und Hardware?}

\textbf{Sensorik:} befasst sich mit der Erfassung physikalischer, chemischer oder biologischer Größen durch Sensoren, die diese in elektrische Signale umwandeln. Diese Signale können anschließend verarbeitet, analysiert und interpretiert werden, um Informationen über die Umgebung oder den Zustand eines Systems zu erhalten.  

\textbf{Hardware:} Umfasst die physischen Komponenten, wie Mikrocontroller, Sensoren und Datenverarbeitungseinheiten, die für die Erfassung und Verarbeitung dieser Signale erforderlich sind.

\subsection{Warum ist Sensorik für die Bewegungserkennung relevant?}

Die präzise Erfassung und Analyse von Bewegungen ist in vielen Bereichen von Bedeutung, beispielsweise in der \emph{Medizintechnik}, \emph{Sportwissenschaft} und \emph{Robotik}. Durch den Einsatz von Sensoren können Bewegungsmuster erkannt, analysiert und klassifiziert werden. Dies ermöglicht unter anderem die Überwachung von körperlichen Aktivitäten, die Steuerung von Prothesen oder die Interaktion mit Robotern.  

\textbf{Mehrfachsensorik:} Bewegungserkennungssysteme nutzen häufig eine Kombination verschiedener Sensoren, um möglichst genaue Daten zu erfassen. Daher ist die Wahl geeigneter Hardware entscheidend.



\section{Auswahl der Hardware}
In diesem Abschnitt wird erläutert, welche Hardware-Komponenten  verwendet werden und warum wir diese verwenden.
\subsection{Arduino Nano 33 BLE REV2}

Für die Erfassung der Sensordaten wird der \texttt{Arduino Nano 33 BLE REV2} verwendet. Dieser Mikrocontroller eignet sich gut für Sensordatenanwendungen, da er eine integrierte \emph{\ac{IMU}} besitzt. Diese Einheit kann \emph{Beschleunigungswerte}, \emph{Drehbewegungen} und \emph{Magnetfelder} erfassen, wodurch er ideal für die Bewegungserkennung ist.  
Ein wesentlicher Vorteil dieses Boards ist seine \emph{\ac{BLE}}-Funktionalität, die eine drahtlose Übertragung der Sensordaten ermöglicht. Dadurch können Messwerte direkt an einen zentralen Rechner gesendet werden, ohne dass eine physische Verbindung erforderlich ist. Wie einem Raspberry Pi

\subsection{Raspberry Pi 4}

Für die Verarbeitung der Sensordaten wird ein \texttt{Raspberry Pi 4} verwendet. Dieser leistungsfähige Einplatinencomputer dient als zentrale Recheneinheit und ist dank seines \textbf{Quad-Core-Prozessors} in der Lage, große Mengen an Sensordaten effizient zu verarbeiten. Seine Schnittstellenvielfalt ermöglicht eine einfache Integration mit dem Arduino und anderen Hardware-Komponenten. Zudem ist der Raspberry Pi für seine breite Unterstützung von Programmiersprachen und Software-Frameworks bekannt, darunter Python, was die Implementierung von Machine-Learning-Modellen zur Bewegungsanalyse erleichtert.

\subsection{Warum diese Kombination?}

Die Kombination aus \texttt{Arduino Nano 33 BLE REV2} und \texttt{Raspberry Pi 4} wurde gewählt, weil beide Geräte ihre jeweiligen Stärken optimal ergänzen:

\begin{itemize}
	\item \textbf{Echtzeit-Datenerfassung:} Der Arduino eignet sich ideal für die direkte Erfassung von Sensordaten, da er direkt mit den Bewegungssensoren verbunden ist und durch seine niedrige Leistungsaufnahme effizient arbeitet.
	\item \textbf{Leistungsstarke Verarbeitung:} Der Raspberry Pi bietet ausreichend Rechenleistung, um die gesammelten Daten zu analysieren, zu speichern und für Machine-Learning-Modelle aufzubereiten.
	\item \textbf{Drahtlose Kommunikation:} Durch die Nutzung von \emph{BLE} können die Sensordaten kabellos an den Raspberry Pi übertragen werden, was eine flexible und skalierbare Architektur ermöglicht.
\end{itemize}

Die folgende Tabelle listet die gängigsten alternativen Hardwareoptionen auf. Diese wurden mit dem Arduino Nano 33 BLE Rev2 und dem Raspberry Pi 4 verglichen. Nachfolgend wird erläutert, warum sie für dieses Projekt nicht geeignet sind:

\bgroup
\def\arraystretch{1.5}
\begin{table}[h!]
	\centering
	\begin{tabular}{|p{3cm}|p{5cm}|p{6cm}|}
		\hline
		\textbf{Hardware- option} & \textbf{Beschreibung} & \textbf{Gründe für die Nichtauswahl} \\
		\hline 
		Arduino Uno & Einfacher Mikrocontroller. & Keine integrierte IMU, was zusätzliche Sensoren erfordert. Keine BLE-Funktionalität für drahtlose Datenübertragung. \\
		\hline
		ESP32 & Mikrocontroller mit höherer Rechenleistung und integriertem Wi-Fi und BLE. & Keine integrierte IMU, was zusätzliche Sensoren erfordert. \\
		\hline
		Raspberry Pi Zero & Kompakter und kostengünstiger Einplatinencomputer. & Weniger Rechenleistung im Vergleich zum Raspberry Pi 4, was die Verarbeitung großer Datenmengen erschwert. \cite{gamess2022performance	}\\
		\hline
		BeagleBone Black & Leistungsfähiger Einplatinencomputer mit vielen Schnittstellen. & Höherer Stromverbrauch und höhere Kosten im Vergleich zum Raspberry Pi 4. Komplexere Einrichtung und weniger Community-Unterstützung. Weniger geeignet für allgemeine Zwecke im Vergleich zum Raspberry Pi. \cite{patil2016comparative} \\
		\hline
	\end{tabular}
	\caption{Vergleich alternativer Hardware-Optionen}
	\label{tab:hardware_vergleich}
\end{table}
\egroup

Aus den in Tabelle \ref{tab:hardware_vergleich} genannten Gründen ist zu erkennen, warum die alternativen Hardwareoptionen Arduino Uno, ESP32, Raspberry Pi Zero und BeagleBone Black für dieses Projekt nicht geeignet sind. Der Arduino Uno und der ESP32 scheiden aufgrund der fehlenden integrierten IMU, die zusätzliche Sensoren erfordern würde, aus. Der Raspberry Pi Zero bietet nicht genügend Rechenleistung für die Verarbeitung großer Datenmengen und der BeagleBone Black ist aufgrund des höheren Stromverbrauchs, der höheren Kosten und des komplexeren Setups weniger geeignet.

Die Wahl fiel daher auf den Arduino Nano 33 BLE REV2 und den Raspberry Pi 4, da diese unseren Annahmen am ehesten entsprachen. Der Arduino Nano 33 BLE REV2 bietet eine integrierte IMU, die Beschleunigungen, Rotationen und Magnetfelder erfassen kann, sowie BLE-Funktionalität für die drahtlose Datenübertragung. Der Raspberry Pi 4 verfügt über ausreichend Rechenleistung, um die gesammelten Daten effizient zu verarbeiten und für Machine-Learning-Modelle aufzubereiten. Diese Kombination ermöglicht eine präzise und flexible Erfassung und Analyse von Bewegungsdaten, was für die Bewegungserkennung entscheidend ist.

Nachdem nun die Hardware-Auswahl erläutert wurde, wird im nächsten Abschnitt auf die verwendeten Sensoren eingegangen und deren Bedeutung für die Bewegungserkennung erläutert.

\section{Sensoren zur Bewegungserkennung}

Um Bewegungen zuverlässig erfassen zu können, werden mehrere Sensortypen kombiniert. Jeder dieser Sensoren liefert spezifische Informationen, die zusammen eine präzise Analyse ermöglichen:

\begin{itemize}
	\item \textbf{Beschleunigungssensor (Accelerometer):} Misst die \emph{lineare Beschleunigung} entlang der drei Raumachsen und ermöglicht die Erkennung von Bewegungen wie Gehen, Laufen oder Springen.
	\item \textbf{Gyroskop:} Erfasst \emph{Drehbewegungen} um die drei Raumachsen und hilft dabei, \emph{Rotationen} und \emph{Orientierungsänderungen} zu bestimmen. In Kombination mit dem Beschleunigungssensor ermöglicht das Gyroskop eine präzise Bestimmung der Bewegungsrichtung und -geschwindigkeit.
	\item \textbf{Magnetometer:} Misst das \emph{Magnetfeld} der Erde und unterstützt die Bestimmung der absoluten Orientierung im Raum. Durch die Kombination von Magnetometerdaten mit denen des Beschleunigungssensors und des Gyroskops kann die Bewegungsrichtung genauer bestimmt und \emph{Driftfehler} korrigiert werden. 
\end{itemize}


